\chapter{Fine-tuning de mBERT B2}
\label{ch:mbert}

Este capítulo corresponde al notebook \texttt{04\_finetuning\_mbert}. Se
describe el reentrenamiento de \texttt{bert-base-multilingual-cased}
~\cite{devlin2019bert} sobre el corpus B2 salvadoreño.

\section{Configuración de entrenamiento}

Los hiperparámetros principales (modo CPU, según \texttt{config/settings.yaml})
se resumen en la Tabla~\ref{tab:hparams}.

\begin{table}[H]
\centering
\caption{Hiperparámetros de fine-tuning B2 (CPU).}
\label{tab:hparams}
\begin{tabular}{lc}
\toprule
Parámetro & Valor \\
\midrule
Modelo base & \texttt{bert-base-multilingual-cased} \\
Directorio de salida & \texttt{models/mbert-sv-b2} \\
Train & 3{,}293 (2{,}144 reales + 1{,}149 sintéticos) \\
Épocas & 4 \\
Batch size (train) & 8 \\
Batch size (eval) & 16 \\
Learning rate & $2 \times 10^{-5}$ \\
Warmup ratio & 0{,}1 \\
Weight decay & 0{,}01 \\
Max length & 128 \\
Label smoothing (B2) & 0{,}05 \\
Early stopping patience & 2 \\
Semilla & 42 \\
\bottomrule
\end{tabular}
\end{table}

\section{Pesos de clase B2}

Los pesos publicados en \texttt{metadata.json} (Tabla~\ref{tab:weights}) mitigan
el desbalance y se combinan con masa efectiva $\approx 90$ en amenazas sintéticas
durante el entrenamiento.

\section{Métricas de validación}

Al finalizar el entrenamiento (época 4), el mejor checkpoint
(\texttt{checkpoint-1644}) alcanza en validación un F1-macro de \textbf{0{,}756}
y una accuracy de \textbf{0{,}756} (\texttt{train\_metrics.json}). La
Figura~\ref{fig:training} resume las métricas finales de validación; el historial
completo de pérdida por step no fue exportado desde Colab.

\begin{figure}[H]
\centering
\includegraphics[width=0.85\textwidth]{09_training_loss.png}
\caption{Curva de entrenamiento: métricas finales de validación B2.}
\label{fig:training}
\end{figure}

El checkpoint oficial se guarda en \texttt{models/mbert-sv-b2/}. Los pesos no
se versionan en Git. El contrato de inferencia (\texttt{inference\_contract.json})
fija \texttt{max\_length=128}, columna \texttt{texto\_modelo} y normalizador
\texttt{normalize\_for\_model} v1.1.

\section{Comparación con el checkpoint v1}

El checkpoint v1 (\texttt{checkpoint-1072}, train 2{,}144) alcanzaba F1-macro
\textbf{0{,}7777} en test. El B2 prioriza robustez adversarial y equidad
dialectal; la evaluación del Capítulo~\ref{ch:evaluacion} cuantifica el
intercambio frente al v1.
